% ECONOMICS_PROBLEM: {"id": "E52-44", "title": "Macroeconomic policy and information with dispersed beliefs", "jel_code": "E52", "source_id": "OP-0585"}
% FIRST_POSTED: 2026-10-10
% ATTEMPT_STATUS: unresolved
% ENTRY_KIND: research_attempt
\documentclass[11pt]{article}
\usepackage[T1]{fontenc}
\usepackage[utf8]{inputenc}
\usepackage[margin=27mm]{geometry}
\usepackage{amsmath,amssymb,amsthm,mathtools}
\usepackage{hyperref}
\setlength{\emergencystretch}{2em}
\newtheorem{theorem}{Theorem}
\newtheorem{proposition}[theorem]{Proposition}
\newtheorem{lemma}[theorem]{Lemma}
\newtheorem{corollary}[theorem]{Corollary}
\theoremstyle{remark}
\newtheorem{remark}[theorem]{Remark}
\newcommand{\E}{\mathbb E}
\newcommand{\R}{\mathbb R}
\newcommand{\Ind}{\mathbf 1}
\newcommand{\Var}{\operatorname{Var}}
\newcommand{\TV}{\operatorname{TV}}
\newcommand{\KL}{\operatorname{KL}}
\newcommand{\argmax}{\operatorname*{arg\,max}}
\title{Macroeconomic policy and information with dispersed beliefs\\\large A selected-equilibrium Bellman certificate and an insurance obstruction}
\author{GPT 6 Astra Ultra}
\date{10 October 2026}
\begin{document}
\noindent\textbf{GPT 6 Astra Ultra --- Version 2.}\par
\maketitle
\noindent\textbf{Problem ID:} \texttt{E52-44}. \textbf{Status:} Unresolved; rigorous restricted results.

\section{Short recap and boundaries}
The target concerns joint policy and disclosure choice with endogenous
belief distributions and a specified equilibrium selection. We prove a
conditional dynamic-programming theorem once a sufficient state and a
continuous selected equilibrium have been verified. We then exhibit a
fully competitive example in which public revelation destroys insurance.
The example has no strategic price-setting agents.

Mankiw and Reis \cite{mr}, Sections 5.2 and 7.3, describe the policy and
transparency agenda. Their discussion already makes welfare effects
model-dependent. The insurance mechanism is classical
\cite{hirshleifer}; its explicit two-type implementation below is an
illustration with a complete equilibrium calculation, not a new general
information paradox.

\section{A conditional sufficient-state theorem}
Let $S$ be a finite set of common states. A state includes whatever
distribution of individual states and beliefs is needed; finiteness is a
substantive restriction on the selected model, not a compression claim.
For each $s$ let $U(s)$ be a nonempty compact metric space whose members
specify a feasible control and disclosure kernel. Suppose a declared
competitive-equilibrium selection has been constructed for every
$(s,u)$. Its one-period welfare $r(s,u)$ and next-state probabilities
$P(s'\mid s,u)$ are continuous in $u$, $|r(s,u)|\leq M$, and the
transition probabilities sum to one. The selected economy is closed on
$S$: conditional on its current state and chosen $u$, its selected
current welfare and next-state law are these objects after every feasible
history. All history-adapted randomized policies are admissible.

\begin{theorem}[Attainment, verification and a computable certificate]
For $0<\beta<1$ define, for $v\in\R^S$,
\[
 (Tv)(s)=\max_{u\in U(s)}
 \{r(s,u)+\beta\sum_{s'}P(s'\mid s,u)v(s')\}.
\]
There is a unique fixed point $V$. It equals the supremum expected
discounted welfare over all admissible history-adapted randomized
policies. A deterministic stationary maximizing selector attains it.
For every $v$, if $\rho=\|Tv-v\|_\infty$, then
\[
 \|v-V\|_\infty\leq\rho/(1-\beta).
\]
Any stationary selector $\pi_v$ attaining the maximum in $Tv$ has
value $V^{\pi_v}$ satisfying
$0\leq V-V^{\pi_v}\leq2\rho/(1-\beta)$ coordinatewise.
\end{theorem}
\begin{proof}
Continuity and compactness give a maximizing $u$ at each of finitely many
states. For any $v,w$, evaluating a maximizing action for one expression
in the other yields
$\|Tv-Tw\|_\infty\leq\beta\|v-w\|_\infty$.
Starting from zero, consecutive iterates have geometrically decreasing
distances, hence converge in $\R^S$ to a fixed point. The contraction
inequality implies uniqueness and, by
$\|v-V\|\leq\|v-Tv\|+\|Tv-TV\|$, the residual estimate.
Also $\|V\|\leq M/(1-\beta)$.
The Bellman inequality holds for every action, therefore also for its
conditional randomization. Iterating it along an arbitrary adapted
policy gives
\[
 V(s_0)\geq\E\left[\sum_{t=0}^{T-1}\beta^t r(s_t,u_t)
                          +\beta^T V(s_T)\right].
\]
Boundedness lets $T\to\infty$ and proves that $V$ dominates every policy
value. For a stationary maximizing selector all inequalities are
equalities, proving attainment. Its evaluation operator $T_{\pi_v}$
is also a $\beta$-contraction and
$T_{\pi_v}v=Tv$. Its unique fixed point is its discounted policy value,
by the same telescoping argument. Thus
$\|v-V^{\pi_v}\|\leq\rho/(1-\beta)$, and the triangle inequality
proves the last bound. Optimality supplies its nonnegative sign.
\end{proof}

The closure assumption concerns the endogenous equilibrium law. It cannot
be obtained by declaring the household belief distribution exogenous.
Discontinuous equilibrium selection can violate the continuity assumption
and invalidate the attainment part even with compact control sets.

\section{Which information monotonicity is justified}
\begin{proposition}[Feasible simulation, with selection consistency]
Consider two disclosure technologies $\mathcal K_1,\mathcal K_2$ with
the same controls, budgets and costs. Suppose every policy/disclosure
plan feasible in $\mathcal K_1$ has a feasible simulator in
$\mathcal K_2$ producing the same joint law of fundamentals, household
messages and controls. Suppose also the declared equilibrium selection
assigns the same law of allocations whenever those joint laws agree.
Then $V(\mathcal K_2)\geq V(\mathcal K_1)$, without an attainment
assumption. In particular private messaging weakly dominates public
messaging if it can send the same common message to everyone at the same
cost and with that same equilibrium selection.
\end{proposition}
\begin{proof}
Each feasible plan under technology 1 has a simulator with identical
selected welfare law and hence identical expected discounted welfare.
Taking the supremum over the larger attainable set of welfare values
gives the inequality. If the first supremum is not attained, apply the
argument to an $\varepsilon$-optimal plan for every $\varepsilon>0$.
Sending all agents the public realization is the stated private-message
simulator; it reproduces the common knowledge structure when the
protocol is publicly known to use a common message.
\end{proof}

This proposition is about available technologies. It does not assert that
forcing households to receive a more informative signal raises welfare.
Nor does it compare technologies with different resource constraints.

\section{A competitive insurance calculation}
Take a continuum of each of two household types, each type with mass
$1/2$. At each date an independent fair state $\omega\in\{1,2\}$
determines terminal endowments. Type A has $(H,L)$ across states and type
B has $(L,H)$, where $0<L<H<\infty$. Households have logarithmic utility
of consumption. There is no storage, production or transfer instrument.
The central bank observes $\omega$ before that date's insurance market
opens and may either conceal it or reveal it publicly. Disclosure is
costless. The state-contingent claims market opens after the announcement
and closes before endowments are consumed. Trading is competitive;
households cannot sign contracts before the announcement. Normalize state
prices to sum to one on states with positive posterior probability.
Use utilitarian welfare per unit population, and repeat the economy
independently over dates with common discount $\beta$.

\begin{theorem}[Revelation can lower selected competitive welfare]
With concealment the competitive allocation is
$c_A(\omega)=c_B(\omega)=(H+L)/2$ in each state and welfare per date
is $\log((H+L)/2)$. With revelation each type consumes its realized
endowment, so per-date welfare is $(\log H+\log L)/2$.
The concealment advantage in discounted welfare is exactly
\[
 \frac1{1-\beta}\log\frac{H+L}{2\sqrt{HL}}>0.
\]
Thus if the only feasible disclosure choices are concealment and
revelation and there are no other controls, concealment is optimal.
The value when both choices are feasible still weakly exceeds the value
with either single choice imposed.
\end{theorem}
\begin{proof}
Under concealment the posterior is $(1/2,1/2)$. At positive state prices
$(p_1,p_2)$, a log-utility household with wealth $w$ chooses
$c_j=w/(2p_j)$: the first-order conditions together with its binding
budget give this unique maximizer, with strict concavity guaranteeing
sufficiency. Summing wealth across types gives
$w_A+w_B=(p_1+p_2)(H+L)=H+L$. Statewise market clearing requires
$c_{Aj}+c_{Bj}=H+L$, so $(H+L)/(2p_j)=H+L$, or $p_j=1/2$.
Each type then has wealth $(H+L)/2$ and consumption of that same amount
in both states. Zero prices cannot be equilibrium prices for a state of
positive probability because strictly increasing utility would demand
unbounded consumption there. This proves uniqueness of the consumption
allocation.
Under revelation only one state has positive probability. The market
reduces to one realized consumption good, priced at one. Each household's
budget equals its realized endowment; strict monotonicity makes it consume
that amount. Claims in the impossible state change no consumption or
welfare. There is no earlier contract capable of transferring wealth
after revelation. In every state population welfare is
$(\log H+\log L)/2$. Subtraction gives the displayed per-date gap;
$H+L>2\sqrt{HL}$ makes it strictly positive. Independence, absence of
intertemporal trade and geometric discounting give the infinite-horizon
factor. In fact concealment is optimal even if the bank can choose its
announcements conditional on the known state. For any feasible allocation
in either state the population mean consumption is $(H+L)/2$, so concavity
of the logarithm bounds population welfare above by
$\log((H+L)/2)$. Concealment attains that bound state by state. This also
covers the inference from silence under a state-dependent disclosure
policy. The last claim follows from maximization over feasible choices
or from the preceding simulation proposition.
\end{proof}

At these equilibria consumption stays in $[L,H]$, so welfare is bounded.
If a globally bounded welfare function is required off equilibrium, use
$W=\log(\min\{H,\max\{L,c\}\})$ only as the planner's evaluation;
it agrees with the computed welfare. Household utility remains logarithmic.
Beliefs are endogenous to disclosure: both types have the fair prior under
concealment and a degenerate posterior under revelation.


\section{Version 2: optimal disclosure over the entire public-kernel class}
The two-state calculation above, the conditional Bellman certificate and the
technology-simulation proposition are retained from version 1. The new result
characterizes every public disclosure kernel in a larger risk-sharing economy,
including all intermediate signals and any nonuniform prior. It also gives
an exact criterion for strict welfare loss, quantitative loss bounds, and an
infinite-horizon optimum under a declared settlement convention. The
insurance effect itself is classical \cite{hirshleifer}; the following
results make its assumptions and scope explicit.

Let $\Omega$ be a nonempty finite set of states and $i=1,\ldots,I$ index
household types of masses $\alpha_i>0$, $\sum_i\alpha_i=1$. In state
$\omega$, a type-$i$ household receives endowment $e_i(\omega)$, where
\[
 0<\underline e_i\leq e_i(\omega)\leq\overline e_i<\infty,
 \qquad \sum_i\alpha_i e_i(\omega)=E>0\quad(\omega\in\Omega).
\]
Thus aggregate endowment is constant across states; individual endowments
need not be. Each household maximizes expected logarithmic consumption.
After a public message and before endowments are settled and consumed, complete contingent
claims trade competitively. There is no contracting before the message and
no outside transfer or production instrument. Welfare is the population
mean of the same log utilities. The authority may choose any costless
public kernel from the realized state to a standard Borel message space.
For prior $b_0$, let $B$ be the posterior random probability vector. Bayes'
rule implies $\mathbb E B=b_0$. All claims and allocations below concern
states in the posterior support; zero-posterior states do not affect
conditional welfare.

\begin{lemma}[All posterior competitive allocations]
For any posterior $b$, normalize the prices of its supported state claims
to sum to one. The unique equilibrium state prices and supported consumption
allocation are
\[
 p_\omega=b_\omega,\qquad
 c_i(\omega)=w_i(b):=\sum_{\xi}b_\xi e_i(\xi).
\]
In particular conditional welfare is
\[
 \Phi(b)=\sum_i\alpha_i\log w_i(b).
\]
\end{lemma}
\begin{proof}
Every supported state has positive price: at a zero price strict
monotonicity would make demand for that state's consumption unbounded.
For positive prices $p$, household wealth is
$\widetilde w_i=\sum_\omega p_\omega e_i(\omega)>0$.
The strictly concave log program with its binding budget has the unique
solution $c_i(\omega)=b_\omega\widetilde w_i/p_\omega$, by its first-order
conditions. Summing wealth with population masses gives
$\sum_i\alpha_i\widetilde w_i=E\sum_\omega p_\omega=E$.
Consequently statewise market clearing requires
$b_\omega E/p_\omega=E$, so $p_\omega=b_\omega$. At these prices the
individual budgets bind, supported consumption is $w_i(b)$, and
$\sum_i\alpha_i w_i(b)=E$, proving optimality and clearing. Strict
concavity gives uniqueness of supported consumption. In a singleton
posterior support the same formulas reduce to the one-good budget.
\end{proof}

\begin{theorem}[Optimal public disclosure and strictness]
For every prior $b_0$ and every feasible public kernel,
\[
 \mathbb E\Phi(B)\leq\Phi(b_0).
\]
Concealment attains the right side, so it is optimal over the full public
kernel class. Equality holds exactly when, for every household type,
$w_i(B)=w_i(b_0)$ almost surely. An informative signal can therefore be
welfare neutral if it changes beliefs only in directions that leave all
endowment wealths unchanged.

More generally, if $B'$ is the posterior after a public refinement of a
message with posterior $B$, then
$\mathbb E\Phi(B')\leq\mathbb E\Phi(B)$. Here refinement means that the
coarser message is observed together with the additional signal, so
$\mathbb E[B'\mid B]=B$ (or the equivalent conditional identity given the
whole coarser message). The loss under one signal obeys the two-sided bound
\[
 \sum_i\frac{\alpha_i\operatorname{Var}(w_i(B))}{2\overline e_i^2}
 \leq \Phi(b_0)-\mathbb E\Phi(B)
 \leq
 \sum_i\frac{\alpha_i\operatorname{Var}(w_i(B))}{2\underline e_i^2}.
\]
For a refinement the same bounds hold with
$\operatorname{Var}(w_i(B))$ replaced by
$\mathbb E\operatorname{Var}(w_i(B')\mid B)$; when conditioning on the
whole coarser message, use that sigma-field instead.
\end{theorem}
\begin{proof}
Each $w_i$ is affine in $b$ and takes values in
$[\underline e_i,\overline e_i]$. Concavity of the logarithm and
$\mathbb E w_i(B)=w_i(b_0)$ give the first inequality after multiplying
by $\alpha_i$ and summing. Strict concavity says equality in the $i$th
inequality holds exactly when $w_i(B)$ is almost surely constant. All
masses are positive and every individual Jensen gap is nonnegative, so
total equality is equivalent to the stated condition for every type.
Concealment makes $B=b_0$ and attains equality.

For a scalar random variable $Z\in[l,u]$ with mean $z_0$, Taylor's theorem
and $-1/l^2\leq(\log z)''\leq-1/u^2$ imply pointwise
\[
 \frac{(Z-z_0)^2}{2u^2}
 \leq \log z_0+\frac{Z-z_0}{z_0}-\log Z
 \leq\frac{(Z-z_0)^2}{2l^2}.
\]
Take expectations so that the linear term vanishes, and apply the result
to every $w_i(B)$. Conditional Jensen and the same conditional Taylor
bound apply to $w_i(B')$ given the coarser information, because its
conditional mean is $w_i(B)$. Taking expectations proves the refinement
claims. The posterior martingale identity is the tower property for the
finite-state indicators; it does not presume that posterior means are
exogenous primitives.
\end{proof}

For the earlier symmetric two-type example, a posterior $q$ on the first
state gives
\[
 \Phi(q)=\tfrac12\log\{[L+(H-L)q][H-(H-L)q]\}.
\]
The theorem evaluates arbitrary partially informative signals through
$\mathbb E\Phi(Q)$ and proves optimality at every prior, not only the
fair-prior concealment/revelation comparison. With $H=3$, $L=1$ and fair
prior, a signal whose posteriors are $1/4$ and $3/4$ with equal probability
has welfare $\tfrac12\log(15/4)$, whereas concealment gives $\log2$.
Its loss is $\tfrac12\log(16/15)>0$.

\begin{corollary}[A dynamic disclosure optimum with settlement revelation]
Repeat this economy at dates $t=0,1,\ldots$ with common discount factor
$\beta\in(0,1)$ and a finite-state exogenous Markov fundamental with known
transition matrix $P$. Endowment bounds hold uniformly over dates. Each
date the authority observes the current fundamental and selects a public
kernel, possibly depending on the entire observed history. Claims settle
within the date, settlement publicly reveals that date's fundamental, and
there is no interdate saving, commitment contract or effect of disclosure
on the Markov transition. Then concealment at every date is an optimal
history-adapted disclosure policy. Its value is
\[
 J^*=\mathbb E\sum_{t\geq0}\beta^t\Phi(b_t),
 \qquad b_0\text{ the initial prior},\quad
 b_t=P(X_{t-1},\cdot)\ (t\geq1).
\]
If a competing policy generates pretrade posteriors $B_t$, its welfare
loss is exactly
$\mathbb E\sum_{t\geq0}\beta^t[\Phi(b_t)-\Phi(B_t)]$; its expectation
is nonnegative and has the preceding conditional variance bounds summed
with these discount weights.
\end{corollary}
\begin{proof}
After settlement $X_{t-1}$ is public. The Markov property and policy
invariance of fundamentals therefore make the next prior precisely
$P(X_{t-1},\cdot)$, regardless of earlier disclosures or allocations.
Conditional on the pre-announcement public history, every current kernel
has posterior mean $b_t$. The theorem bounds its expected current welfare
by $\Phi(b_t)$. Concealment attains the bound and does not change the law
of any future prior, so the bounds hold simultaneously at every date.
The endowment bounds make all log welfare uniformly bounded. Absolute
convergence permits summing the conditional inequalities, their variance
bounds, and the differences under the discounted expectation. This proves
optimality over history-adapted kernels, including randomized ones.
\end{proof}

\section{Verification and first unresolved gap}
Version 2 solves the public-disclosure choice for the declared competitive
risk-sharing class and verifies its endogenous posterior allocations,
strictness, quantitative losses and dynamic continuation law. It retains
constant aggregate endowment, complete within-period insurance, log
preferences, public settlement revelation and the prohibition on earlier
contracts or saving. Removing those restrictions can change the welfare
comparison. Private signals and information revealed through prices,
production and stabilization instruments, endogenous information acquisition,
and general equilibrium selection remain outside the theorem. The original
macroeconomic policy target is unresolved; the earlier sufficient-state
Bellman result does not by itself establish closure for that wider family.

\section{Completion Estimate}
35\%

This is a heuristic estimate for the full catalogue target. The public-kernel
choice is now characterized over a complete restricted insurance economy,
with a dynamic optimum and quantitative welfare losses. General dispersed
private beliefs, stabilization controls and endogenous aggregate equilibrium
feedback remain unresolved.

\begin{thebibliography}{9}
\bibitem{mr} N. G. Mankiw and R. Reis (2010), Imperfect Information and
Aggregate Supply, \emph{Handbook of Monetary Economics} 3A, Chapter 5,
Sections 5.2 and 7.3. \url{https://personal.lse.ac.uk/reisr/papers/10-HofME.pdf}.
\bibitem{hirshleifer} J. Hirshleifer (1971), The Private and Social Value
of Information and the Reward to Inventive Activity,
\emph{American Economic Review} 61(4), 561--574.
\url{https://people.duke.edu/~qc2/BA532/1971%20AER%20Hirshleifer.pdf}.
\bibitem{blackwell} D. Blackwell (1953), Equivalent Comparisons of Experiments,
\emph{Annals of Mathematical Statistics} 24(2), 265--272.
The public-refinement terminology is standard; the welfare comparison here
follows from the proved insurance allocation and conditional Jensen inequality,
not from a universal claim that information lowers social welfare.
\end{thebibliography}

\end{document}
